\section{The canonical braided unitary tensor embedding}
\label{sec:comparison}

The preceding sections provide the analytic and positivity inputs for
the categorical comparison.  We now identify the tensor structure of
the actual CWX functor.  Essential surjectivity will be proved in the
next section; at this stage we obtain an embedding and finite index
for every sector in its image.

Here $\operatorname{Rep}(A_V)$ denotes the category of compatible
families of locally normal representations of the interval algebras,
with bounded operators intertwining every local action as morphisms.
On the separable Hilbert completions considered here, these are the
transportable DHR representations of the circle net; see
\cite[Section~6]{KLM}.

\begin{theorem}[CWX tensor embedding]
\label{thm:cwx-embedding}
Let $V$ be simple, unitary and strongly rational.  The Hilbert
completion and the local actions of Theorem~\ref{thm:integration}
define a fully faithful $*$-functor
\[
 F^V_{\mathrm{CWX}}:
 \operatorname{Rep}^u(V)\longrightarrow\operatorname{Rep}(A_V).
\]
It has a unitary tensor structure identifying the Huang--Lepowsky
fusion and braiding with Connes fusion and DHR braiding, respectively.
In particular every object in its image is dualizable, hence has
finite index.
\end{theorem}

\begin{proof}
We first verify the analytic inputs to Gui's comparison and identify
the underlying functor. The additional localizability conditions will be
checked before invoking the tensor comparison.

\smallskip
\noindent\emph{Generating fields and charges.}
We verify the hypotheses of Gui's comparison theorem
\cite[Section~3.4 and Corollary~3.23]{GuiUnbounded}.
Regularity follows from rationality and $C_2$-cofiniteness.
Theorem~\ref{thm:strong-unitarity} supplies positive invariant forms
on all simple modules; finite orthogonal sums then supply strong
unitarity on the grading-restricted category.
Theorem~\ref{thm:energy} gives energy bounds on every unitary module,
and Corollary~\ref{cor:locality} gives strong locality of $V$.

Take a generating family $E$ of homogeneous quasi-primary vectors in
$V$ as in the energy argument.  Since strong locality holds for all
vacuum fields, it holds in particular for $E$.
Choose one representative of every simple module, equipped with its
positive form, as the family of generating charges.  This family
certainly tensor-generates the algebraic category.  In each charge
module choose a nonzero lowest-weight vector $a$.  It is quasi-primary
because $L_1$ lowers conformal energy.
Corollary~\ref{cor:charged-general} and
Lemma~\ref{lem:charged-tools} give polynomial energy bounds for every
intertwining operator between unitary modules with this charge.
Proposition~\ref{prop:strong-intertwining} gives the strong
intertwining property for every homogeneous observable, hence for $E$.
The same fixed vector $a$ works for every irreducible intertwining
operator with that charge; it has not been chosen separately for
different sources or targets.

These are exactly Condition~A of \cite[Section~3.4]{GuiUnbounded}.
The chosen charge vectors are quasi-primary, so Condition~B holds as
well.  Thus \cite[Theorem~3.15]{GuiUnbounded} identifies the canonical
positive fusion structures with a unitary modular tensor category.
This agrees with the normalization of
Theorem~\ref{thm:positive-fusion}; positivity of an unspecified
sesquilinear form is not being substituted for canonical positivity.

\smallskip
\noindent\emph{The underlying fully faithful functor.}
The functor on objects is the representation characterized by
\eqref{eq:exact-cwx-smear}.  On a module morphism it is its bounded
extension to Hilbert completions.  Such extensions are bounded because
the modules are finite orthogonal sums of simple modules, and a map
between simple unitary modules is either zero or a scalar multiple
of a unitary isomorphism.  The Hilbert adjoint is again a module map.
The exact smear identity shows that these extensions intertwine the
local algebra actions.  They are therefore the morphism maps of the
CWX functor, rather than of an independently chosen equivalence.
Full faithfulness, including preservation of adjoints, is
\cite[Theorem~3.17]{GuiUnbounded}; see also
\cite[Theorem~4.24]{GuiExtensions}.  The uniqueness in
Theorem~\ref{thm:integration} identifies the functor there with ours.
This full-faithfulness result uses the observable-smear argument of
\cite[Theorem~4.3]{GuiUnitaryII}: localized stress-energy smears recover
$L_0$ and then the algebraic module maps. It precedes the construction
of a categorical extension and does not use tensor comparison.

\smallskip
\noindent\emph{Localizability of finite sums and products.}
We use the corrected localizability requirements in
\cite[item~3]{GuiUnboundedErrata}. Here is their verification for
the operators entering Gui's categorical partial extension.
Write $D_j=1+L_0$ on the Hilbert completion of a unitary module.
The charged energy bounds and the $L_0$-commutator formula give, for
each smeared field $T:\mathcal H_j^\infty\to\mathcal H_k^\infty$
and each nonnegative integer $s$, an integer $r_s$ and a constant $C_s$
such that
\[
 \|D_k^sT\xi\|\leq C_s\|D_j^{r_s}\xi\|.
\]
Indeed, commuting $L_0$ past a smear differentiates its smooth test
and adds its fixed conformal-weight term; iteration reduces the
estimate to the polynomial mode bounds. The same argument applies
to the formal adjoint. This is the Sobolev estimate used in
\cite[Section~3.1 and Proposition~3.2]{GuiUnbounded}.

Every bounded net morphism $S$ between objects of the CWX image is,
by the full faithfulness just established, the bounded extension of
a VOA morphism. All intermediate objects here are completions of VOA
fusion products, on which the local actions are already defined. Thus
$S$ and $S^*$ commute with the respective $L_0$ operators and preserve
all Sobolev spaces. Gui's left and right operators are products of a
charged smear with finitely many observable smears
\cite[Equations~(3.18)--(3.19)]{GuiUnbounded}. Their finite compositions,
and those compositions followed by a bounded net morphism, therefore
have the form covered by the preceding estimates. Consequently each
finite sum, with common source and target,
\[
 B=\sum_{\alpha=1}^{m}c_\alpha S_\alpha
       T_{\alpha,n_\alpha}\cdots T_{\alpha,1}
       \big|_{\mathcal H_j^\infty}
\]
is smooth, has a densely defined formal adjoint, and satisfies
$\|B\xi\|\leq C\|D_j^r\xi\|$ for some integer $r$.
In particular $B$ is closable. This includes the finite sums of left
or right fields in the corrected Definition~2.29 and the finite
sums with a bounded morphism required by corrected Definition~2.41.

To verify localizability, let $\mathscr D_0\subset\mathcal H_j^\infty$
be a dense linear subspace that is quasi-rotation invariant: for some $\delta>0$ there is
a dense $\mathscr D_\delta\subset\mathscr D_0$ such that
$e^{itL_0}\mathscr D_\delta\subset\mathscr D_0$ for $|t|<\delta$.
Then $\mathscr D_0$ is a core for $D_j^r$
\cite[Lemma~7.2]{CKLW}. To see the regularization explicitly, choose
$\varphi\in C_c^\infty(-\delta,\delta)$ and set
\[
 R_\varphi=\int\varphi(t)e^{itL_0}\,dt.
\]
The rapid decay of the Fourier transform of $\varphi$ makes
$D_j^rR_\varphi$ bounded. Hence $R_\varphi$ turns Hilbert-norm
approximations from $\mathscr D_\delta$ into Sobolev-norm
approximations. For a vector in $\mathscr D_\delta$, its convolution
is itself a Sobolev limit of Riemann sums in $\mathscr D_0$.
Taking smooth approximate identities for $\varphi$ now approximates
every vector of $\operatorname{Dom}D_j^r$ in that graph norm.
The displayed bound for $B$ turns this approximation into graph
approximation for $B$ on $\mathcal H_j^\infty$. Hence
$\overline{B|_{\mathscr D_0}}=\overline B$, as required.
Thus the strengthened finite-sum localizability hypotheses hold.

\smallskip
\noindent\emph{Connes fusion, braiding, and finite index.}
Full faithfulness alone would only allow us to transport the VOA
tensor structure onto its image.  The additional comparison with the
already defined net tensor structure is
\cite[Corollary~3.23]{GuiUnbounded}.
It supplies unitary natural maps
\[
 U_{M,N}:F^V_{\mathrm{CWX}}(M)
       \boxtimes_{A_V}F^V_{\mathrm{CWX}}(N)
       \longrightarrow F^V_{\mathrm{CWX}}(M\boxtimes N),
\]
compatible with the unit and associativity constraints.  Here
$\boxtimes_{A_V}$ denotes Connes fusion, while $\boxtimes$ denotes VOA
fusion.  Denote the VOA and net braidings by $c^V$ and
$c^{\mathrm{DHR}}$, using the same circle orientation.
The maps satisfy
\[
 F^V_{\mathrm{CWX}}(c^{V}_{M,N})\,U_{M,N}
   =U_{N,M}\,
      c^{\mathrm{DHR}}_{F^V_{\mathrm{CWX}}(M),F^V_{\mathrm{CWX}}(N)}.
\]
Thus the comparison is braided and unitary for the specified fusion
and braiding on the net side.

Finally, the source category is rigid.  Applying this unitary tensor
functor to its evaluation and coevaluation maps preserves the
conjugate equations.  Their images exhibit a dual for every object
in the image.  The equivalence between dualizability and finite index
for the net sectors then gives the last assertion.
\end{proof}

The image is closed under finite sums and subobjects.  For subobjects,
full faithfulness lifts a net projection to a self-adjoint idempotent
module endomorphism; its range is a unitary submodule, whose completion
is the prescribed subrepresentation.  This proves closure of the
image, but does not yet rule out a simple sector outside it.
The next section establishes precisely that exhaustion statement.
